\begin{titlepage}
\thispagestyle{empty}
\begin{tikzpicture}[remember picture, overlay]
  % Верхняя цветная полоса
  \fill[accent] (current page.north west) rectangle ([yshift=-10.6cm]current page.north east);

  % Декоративные «контейнеры»
  \begin{scope}[shift={([xshift=-6.3cm,yshift=-3.1cm]current page.north east)}]
    \foreach \dx/\dy/\op in {0/0/0.20, 1.15/0.55/0.14, 2.3/1.1/0.08}{
      \draw[white, opacity=\op, fill=white, fill opacity=0.035, line width=0.7pt, rounded corners=1.5pt]
        (\dx,\dy) rectangle ++(1.9,1.55);
      \draw[white, opacity=\op, line width=0.7pt]
        (\dx+0.35,\dy+1.55) -- (\dx+0.35,\dy+0.55) -- (\dx+1.9,\dy+0.55);
    }
  \end{scope}

  \node[anchor=north west, text width=12cm, align=left, font=\sffamily]
    at ([xshift=2.7cm,yshift=-2.7cm]current page.north west) {%
      {\color{white!80!accent}\bfseries\small\MakeUppercase{Учебное пособие}}\\[0.75cm]
      {\color{white}\bfseries\fontsize{52}{56}\selectfont ОСНОВЫ\\[3pt]DOCKER}\\[0.65cm]
      {\color{white!90!accent}\Large Контейнеры, образы,\\ сети и Docker Compose}
  };
\end{tikzpicture}

\vspace*{11.4cm}

\begin{center}
\begin{tikzpicture}[font=\sffamily\footnotesize, x=1cm, y=1cm]
  \draw[fill=white, draw=warm, dashed, rounded corners=2pt, line width=0.8pt]
    (0,2.85) rectangle (12.5,3.7);
  \node[anchor=west, text=warm] at (0.3,3.275) {записываемый слой контейнера};

  \draw[fill=accent!16, draw=accent!55!black, rounded corners=2pt]
    (0,1.9) rectangle (12.5,2.75);
  \node[anchor=west, text=accent!80!black] at (0.3,2.325)
    {слой~3: \texttt{COPY . .} --- код приложения};

  \draw[fill=accent!10, draw=accent!55!black, rounded corners=2pt]
    (0,0.95) rectangle (12.5,1.8);
  \node[anchor=west, text=accent!80!black] at (0.3,1.375)
    {слой~2: \texttt{RUN npm install} --- зависимости};

  \draw[fill=accent!6, draw=accent!55!black, rounded corners=2pt]
    (0,0) rectangle (12.5,0.85);
  \node[anchor=west, text=accent!80!black] at (0.3,0.425)
    {слой~1: \texttt{FROM node:22-alpine} --- базовый образ};

  \node[anchor=north, text=muted, align=center] at (6.25,-0.25)
    {Образ~--- неизменяемый стек слоёв. Контейнер добавляет сверху\\один записываемый слой, который исчезает вместе с контейнером.};
\end{tikzpicture}
\end{center}

\vfill

\begin{center}
  {\sffamily\small\color{muted}
   Четыре части \textbullet{} 13~глав\\
   Схемы, листинги и вопросы для самопроверки}\\[0.5cm]
  {\sffamily\bfseries\large\color{accent} 2026}
\end{center}

\end{titlepage}
